\clearpage
\fsection{Arquitectura del sistema}

Desde el punto de vista técnico, el sistema está compuesto por los siguientes elementos interconectados mediante
protocolo Profinet en la red 192.168.15.xxx:
\begin{itemize}
	\item Cámara Orbbec Femto Mega (192.168.15.41): cámara de profundidad 3D que captura tanto imagen en
	      color como nube de puntos, proporcionando la información espacial necesaria para localizar
	      las piezas en el contenedor.
	\item PC Industrial Siemens SIMATIC IPC BX-59A (192.168.15.19): núcleo de procesamiento del sistema.
	      Ejecuta tres contenedores Docker: uno con los drivers de la cámara, otro con el frontend
	      web de configuración y monitorización, y un tercero con el modelo de inteligencia artificial
	      preentrenado por Siemens para la detección del contenedor y las piezas, así como el cálculo
	      del punto de agarre.
	\item PLC Siemens S7-1500, modelo 1516-3 PN/DP (192.168.15.10): controlador industrial que gestiona
	      la lógica del proceso. Se comunica con el IPC para solicitar coordenadas de agarre y con
	      el robot para enviar las posiciones de destino. Programado mediante TIA Portal v18.
	\item Robot colaborativo Universal Robots UR3e (192.168.15.20): brazo robótico de 6 ejes que
	      ejecuta las operaciones de picking. Recibe las coordenadas calculadas por la IA y
	      realiza el movimiento de agarre y depósito de la pieza.
	\item Pinza de succión On-Robot: herramienta de efector final montada en la muñeca del robot,
	      que realiza el agarre de las piezas mediante vacío.
\end{itemize}

